%==============================================================================
% Section 5 -- 2D Projection and EWA Splatting
%==============================================================================
\section{2D Projection and Jacobian Approximation}
\label{sec:projection}

\subsection{Overview}

Each 3D Gaussian must be projected onto the image plane. The key result is that
the projection of a 3D Gaussian under a perspective camera yields
\emph{approximately} a 2D Gaussian, following the \textbf{EWA splatting}
framework~\cite{zwicker2001ewa}.

%------------------------------------------------------------------------------
\subsection{Camera Model and Mean Projection}

A world-space point maps to camera space as
$\mathbf{x}_c = \mathbf{R}_{cw}\,\mathbf{x}_w + \mathbf{t}_{cw}$,
then projects to image coordinates:
\begin{equation}
  \pi(\mathbf{t}) =
  \begin{pmatrix} f_x\,t_x/t_z + c_x \\ f_y\,t_y/t_z + c_y \end{pmatrix}
\end{equation}
The 2D projected mean is simply $\mean'_i = \pi(\mathbf{R}_{cw}\,\mean_i +
\mathbf{t}_{cw})$.

%------------------------------------------------------------------------------
\subsection{Covariance Projection via Jacobian Approximation}

Since $\pi$ is nonlinear, a first-order Taylor expansion is used. The Jacobian
of $\pi$ at camera-space point $\mathbf{t}_i = (t_{i,x}, t_{i,y}, t_{i,z})$ is:
\begin{equation}
  \mathbf{J}_i =
  \begin{pmatrix}
    f_x/t_{i,z} & 0 & -f_x\,t_{i,x}/t_{i,z}^2 \\
    0 & f_y/t_{i,z} & -f_y\,t_{i,y}/t_{i,z}^2
  \end{pmatrix}
  \in \R^{2\times 3}
\end{equation}

The \textbf{projected 2D covariance} is:
\begin{equation}
  \boxed{
    \cov'_i = \mathbf{J}_i\,\mathbf{R}_{cw}\,\cov_i\,\mathbf{R}_{cw}\transp\,\mathbf{J}_i\transp
    \in \R^{2\times 2}
  }
  \label{eq:cov_proj}
\end{equation}

The resulting 2D splat in screen space is:
\begin{equation}
  \Gcal'_i(\mathbf{u}) =
    \exp\!\left(
      -\tfrac{1}{2}(\mathbf{u} - \mean'_i)\transp(\cov'_i)^{-1}(\mathbf{u} - \mean'_i)
    \right), \quad \mathbf{u} \in \R^2
\end{equation}
$(\cov'_i)^{-1}$ is computed in closed form since it is $2\times2$.

%------------------------------------------------------------------------------
\subsection{Bounding Box and Tile Culling}

For each splat, a bounding box of half-extent $r = 3\sqrt{\lambda_{\max}(\cov'_i)}$
(the $3\sigma$ ellipse, capturing $>99\%$ of the Gaussian's mass) determines
which screen tiles overlap with the splat. Pixels outside this radius are
discarded, reducing per-tile Gaussian counts significantly.

The Jacobian-based approach is chosen because it is accurate for typical scene
scales and fully differentiable, enabling gradient flow through the projection.
